\begingroup
\renewcommand{\arraystretch}{1.45}
\Needspace{42\baselineskip}
\InterpolationArticle{29}
There remain the cases in which the function $X$ contains only cosines or only sines, which we treated generally in Articles 16, 17. As before, we shall suppose that a period of $\mu$ terms is given, say $X=A,B,C,D$ etc. for $x=a,b,c,d$ etc. Two cases must, however, be carefully distinguished here, since they require different treatment:

the $1^{\mathrm{st}}$ case, in which $a$ is the complement of one of the following values $b,c,d$ etc. to $360^{0}$ or to $0$ or to a multiple of the circumference. Let this term be $a+\frac{k}{\mu}360^{0}$,
\SourceLeaf{320}
and let $a+\frac{k}{\mu}360^{0}+a=l\times360^{0}$, or $a=(\mu l-k)\frac{180}{\mu}$. Hence, in this case, $a$ will be divisible by $\frac{180}{\mu}$, and therefore $\sin\mu a=0$; furthermore, we have $b=(\mu l-k+2)\frac{180^{0}}{\mu}$, and $\sin\mu b=0$, and similarly for the following values $c,d$ etc. Thus, supposing $a$ to be less than $\frac{360^{0}}{\mu}$ but nonnegative (which is permissible, since we may subtract from each of $a,b,c,d$ etc. the nearest smaller multiple of the whole circumference and take the smallest of the remainders, which will have that property, as the beginning of the period), there will be either $a=0$, and $b=\frac1\mu360^{0}$, $c=\frac2\mu360^{0}$, $d=\frac3\mu360^{0}$ etc.; or $a=\frac1\mu180^{0}$, $b=\frac3\mu180^{0}$, $c=\frac5\mu180^{0}$, $d=\frac7\mu180^{0}$ etc.

\Needspace{22\baselineskip}
and the $2^{\mathrm{nd}}$ case, in which it is not the complement of any of the following values $b,c,d$ etc. to $360^{0}$ or $0$ or a multiple of the circumference. Here $a$ will not be divisible by $\frac{180^{0}}{\mu}$, and therefore $\sin\mu a=\sin\mu b=\sin\mu c=\sin\mu d$ etc. is not $=0$.

In the latter case, the general method of Article 16 can be applied immediately. In the former, however, from the values $a,b,c,d$ etc. (which we suppose reduced below $360^{0}$), all those greater than $180^{0}$ must first be discarded, since their complements to $360^{0}$ occur among the remaining values; and, in addition, the values $0$ and $180^{0}$, when present, if the function $X$ contains only sines. But since applying the general method in this way would be less convenient, we shall handle this case differently below. We shall now begin with the case in which $\sin\mu a$ is not $=0$.
\Needspace{37\baselineskip}
\InterpolationArticle{30}
First let $X$ be a function of the form
\[
\begin{aligned}
&\alpha+\alpha'\cos x+\alpha''\cos2x+\alpha'''\cos3x+.\,.\,.+\alpha^{n}\cos nx
\end{aligned}
\]
so that $\mu=n+1$ must hold. Here, by the second lemma of Article 19, we have
\[
\begin{aligned}
&(\cos t-\cos b)(\cos t-\cos c)(\cos t-\cos d)\ \text{etc.}
\\[5pt]
&=\frac1{2^{\mu-1}}\times\frac{\cos\mu t-\cos\mu a}{\cos t-\cos a}\\[5pt]
&=\frac1{2^{\mu-1}\sin a}\{\sin\mu a+2\sin(\mu-1)a\cos t\\[5pt]
&\qquad{}+2\sin(\mu-2)a\cos2t+\text{etc.}+2\sin a\cos(\mu-1)t\}
\end{aligned}
\]
Indeed, it is easily confirmed that multiplying the sum in the second part of this equation by $\cos t-\cos a$ gives $=\sin a(\cos\mu t-\cos\mu a)$. Hence, on writing $a$ for $t$, there also results
\[
\begin{aligned}
&(\cos a-\cos b)(\cos a-\cos c)(\cos a-\cos d)\ \text{etc.}
\\[5pt]
&=\frac1{2^{\mu-1}\sin a}\{\sin\mu a+2\sin(\mu-1)a\cos a\\[5pt]
&\qquad+2\sin(\mu-2)a\cos2a+\text{etc.}+2\sin a\cos(\mu-1)a\}
\end{aligned}
\]
\Needspace{14\baselineskip}
The sum in this expression is $=\mu\sin\mu a$, as is evident from the fact that
\SourceLeaf{321}
\[
\begin{aligned}
&\text{second term}=\sin\mu a+\sin(\mu-2)a\\[5pt]
&\text{last term}=\sin\mu a-\sin(\mu-2)a\\[5pt]
&\text{further, third term}=\sin\mu a+\sin(\mu-4)a\\[5pt]
&\text{next-to-last term}=\sin\mu a-\sin(\mu-4)a\\[5pt]
&\text{etc.}
\end{aligned}
\]
\Needspace{44\baselineskip}
It follows that the coefficient of $A$ in the formula of Article 16 for $T$ becomes
\[
\begin{aligned}
&\frac1{\mu\sin\mu a}\{\sin\mu a+2\sin(\mu-1)a\cos t\\[5pt]
&\qquad{}+2\sin(\mu-2)a\cos2t+\text{etc.}+2\sin a\sin(\mu-1)t\}
\end{aligned}
\]
Exactly similar expressions result for the coefficients of $B,C,D$ etc., merely changing $a$ into $b,c,d$ etc. Thus, on setting
\[
\begin{aligned}
&\varepsilon=\frac1\mu(A+B+C+D+\text{etc.})\\[5pt]
&\varepsilon'=\frac2{\mu\sin\mu a}\{A\sin(\mu-1)a+B\sin(\mu-1)b\\[5pt]
&\qquad{}+C\sin(\mu-1)c+D\sin(\mu-1)d+\text{etc.}\}\\[5pt]
&\varepsilon''=\frac2{\mu\sin\mu a}\{A\sin(\mu-2)a+B\sin(\mu-2)b\\[5pt]
&\qquad{}+C\sin(\mu-2)c+D\sin(\mu-2)d+\text{etc.}\}\\[5pt]
&\varepsilon'''=\frac2{\mu\sin\mu a}\{A\sin(\mu-3)a+B\sin(\mu-3)b\\[5pt]
&\qquad{}+C\sin(\mu-3)c+D\sin(\mu-3)d+\text{etc.}\}\\[5pt]
&\text{etc.}\\[5pt]
&\varepsilon^{n}=\frac2{\mu\sin\mu a}(A\sin a+B\sin b\\[5pt]
&\qquad{}+C\sin c+D\sin d+\text{etc.})
\end{aligned}
\]
there will be
\[
\begin{aligned}
&T=\varepsilon+\varepsilon'\cos t+\varepsilon''\cos2t+\varepsilon'''\cos3t+\text{etc.}\\[5pt]
&\qquad{}+\varepsilon^{n}\cos nt
\end{aligned}
\]
Since this formula must hold generally for any value of $t$, it will necessarily become identical with $X$ on changing $t$ into $x$, and therefore the coefficients $\varepsilon,\varepsilon',\varepsilon''$ etc. will be equal respectively to $\alpha',\alpha',\alpha''$ etc.
\endgroup
